\chapter{Lista Completa de Componentes}
\label{anx:a}

\orientacao{Os anexos guardam o material que interrompe a leitura mas precisa ficar registrado. Um anexo com o inventário do que foi entregue evita a discussão sobre o que foi entregue. Cada anexo é citado no corpo do relatório. Remova os anexos que o projeto não tem.}

\ph{Lista completa de componentes, com especificação e fornecedor.}

\chapter{Estrutura e Dicionário dos Dados}
\label{anx:b}

\ph{Estrutura de diretórios dos conjuntos de dados entregues e dicionário de dados.}

\chapter{Repositórios e Versões Entregues}
\label{anx:c}

\ph{Endereço de cada repositório e a versão entregue (\emph{commit} ou \emph{tag}). Não cole trecho extenso de código: aponte o repositório. Registre a licença de bibliotecas, conjuntos de dados e modelos de terceiros quando ela impuser condição ao uso pelo parceiro.}

\chapter{Publicações e Divulgação}
\label{anx:d}

\ph{Publicações e material de divulgação produzidos no projeto, com a menção ao apoio da EMBRAPII.}

\chapter{Registros de Propriedade Intelectual}
\label{anx:e}

\ph{O que foi depositado no INPI e em que estágio está. Consulte o Núcleo de Inovação antes de descrever em detalhe qualquer resultado com potencial de proteção.}

\chapter{Material de Capacitação}
\label{anx:f}

\ph{Material de capacitação entregue ao parceiro.}
